\chapter{Before Adding the Board to QEMU}

Adding a new board to QEMU involves multiple steps. This chapter outlines the detailed process we followed to integrate the \textit{S32K3X8EVB board}, based on \textit{S32K358 SoC}, into QEMU. Our implementation leveraged QEMU’s modular architecture, reusing pre-existing CPU models (Cortex-M7) while defining the board, SoC, peripherals, and memory mappings required to emulate this embedded platform. 

\section{What is a System-on-Chip?}
\begin{center}
    Referred as SoC.
\end{center}
A System-on-Chip integrates all major components of a computing system, such as the CPU, memory, I/O peripherals, and other essential controllers, into a single silicon chip. In embedded systems, the SoC is the heart of the hardware platform, providing the necessary computational and I/O capabilities tailored for specific tasks. 

\vspace{0.2cm}

\noindent For example:
\begin{itemize}
    \item The \textit{CPU} (Central Processing Unit) handles the computational logic.
    \item Integrated \textit{peripherals} like UARTs, timers, and memory controllers enable communication, timekeeping, and data storage.
\end{itemize}

\vspace{0.2cm}

\noindent For the project, the System on Chip (SoC) used is the \textit{S32K358 SoC}. It comes pre-configured with a Cortex-M7 CPU and includes additional peripherals such as UART and timers.

\section{Relationship between IRQs and Peripherals}
\begin{center}
    IRQs (Interrupt Requests) are signals generated by peripherals to notify the CPU of events that require immediate attention.
\end{center}
Timers and UARTs are examples of peripherals that generate IRQs. A timer is able to generate IRQs to signal the end of a countdown or periodic event. For example, A Periodic Interrupt Timer (PIT) signals the CPU for periodic tasks, such as updating a clock or managing real-time scheduling. Similarly, a UART generates IRQs to signal the arrival of new data or the completion of a transmission. For example, when a UART receives data from a serial line, it triggers an IRQ to notify the CPU to read the data buffer.

\vspace{0.5cm}

\noindent Each peripheral is assigned a specific IRQ line (unique lines connecting the peripheral to the CPU or interrupt controller) and an interrupt vector (a unique identifier in the vector table pointing to the Interrupt Service Routine for that peripheral). 